\begin{proof}
\begin{enumerate}[label=\arabic*.]
\item On the punctured base the evaluation map is
\[
 E_c:\Lambda_{\mathbb R}\longrightarrow F_c^\vee,
 \qquad E_c(x)(\omega)=\omega(x).
\]
Its kernel is zero: a real vector annihilating $F_c$ also annihilates
$\overline{F_c}$, and these two spaces span the complexified cohomology
space. Both real dimensions are six. If $\Omega$ is the evaluation
matrix of a local frame, its inverse is explicitly
\[
 E_c^{-1}(z)=R^{-1}
       \begin{pmatrix}\operatorname{Re}z\\\operatorname{Im}z\end{pmatrix},
 \qquad
 R=\begin{pmatrix}\operatorname{Re}\Omega\\\operatorname{Im}\Omega\end{pmatrix}.
\]
Under a frame change $\Omega'=G\Omega$, both $R$ and the real coordinate
column of $z$ are multiplied by the real matrix of $G$. Thus this
inverse glues. Evaluation sends each element of $\Lambda$ to its
specified period section and induces the marked real torus bundle
isomorphism, preserving the zero section. For a peripheral transport
$\gamma$ with matrix $M$, the actual evaluation square is
\[
 E_{c,\gamma p}(Mx)=\Phi_\gamma E_{c,p}(x),
\]
where $\Phi_\gamma$ is the transport in the period quotient. Consequently
\[
 J_c=E_c^{-1}iE_c=R^{-1}J_{\rm std}R,
 \qquad
 J_{\rm std}=\begin{pmatrix}0&-I_3\\I_3&0\end{pmatrix},
 \qquad J_c^2=-I_6,
\]
and $J_c(\gamma p)=M J_c(p)M^{-1}$.

\item At a finite site put $u=r e^{2\pi ia}$ on the universal angular
cover and $L=\log r/(2\pi)$. The prescribed attaching translation in
Theorem~\ref{night:P1.10} becomes
\[
 E_c^{-1}\!\left(-\frac{\log u}{2\pi i}E_c(v)\right)
       =-av+LJ_c(u)v.
\]
The affine action is $(x,u)\mapsto(Mx+v/m,\zeta_m u)$, and $Mv=v$ by
Theorem~\ref{night:P1.1}. Write
\[
 \alpha_0(x,u)=x-av,
 \qquad h_c(u)=LJ_c(u)v,
 \qquad \alpha_c=\alpha_0+h_c.
\]
The angular shift is $a\mapsto a+1/m$, and the displayed transport of
$J_c$ gives the exact identities
\[
 h_c(\zeta_m u)=M h_c(u),\qquad
 \alpha_c(Mx+v/m,\zeta_m u)=M\alpha_c(x,u).
\]
Hence $h_c$ is a section of the real local-system vector bundle on
the attaching annulus. Choose a scalar cutoff $\chi$ which is one at
the inner boundary of an attaching collar and zero at its outer
boundary. Translation by $-\chi h_c$ on that collar has inverse
translation by $\chi h_c$; the translations by $-s\chi h_c$,
$0\leq s\leq1$, give an equivariant isotopy. At the inner boundary
it carries $\alpha_c$ to $\alpha_0$ and at the outer boundary it is
the identity. The two finite collars are disjoint.

These maps identify the finite attaching diagrams with the marked
finite-cover diagrams of Theorem~\ref{maps:src:F.2.6}. In particular
they carry the primitive boundary loops $(g^m h^{-1},g)$, their
subdivided mapping prisms and their integral restriction maps to
the ones in that proof. The induced fibre cohomology map of the
translation is the identity through the specified isotopy. Thus the
three order-four torsion transgressions retain the unit coefficients
computed there.

\item The cusp cap is the one constructed in
Theorems~\ref{night:P1.9} and~\ref{night:P1.10}, in its specified
$S$ marking. The actual collapse of Theorem~\ref{night:P1.11} has
primitive image and saturated kernel. Its dual restriction image is
\[
 J_{0,q}=\ker\!\left(\bigwedge^q M_0^t-I\right)
          \subseteq\bigwedge^q\Lambda^\vee;
\]
its additional free kernel has rank one in degree four and is zero
in degrees one through three. This determines the cusp restriction
homomorphism in the marked exterior lattices.

Together with Steps~1--2, the identity maps of those marked lattices
identify the whole restriction diagram with the one calculated in
Theorem~\ref{maps:src:F.2.6}. In particular the cusp column of the
integral complex is $(-j_0,-A_qj_0)$, where
$A_q=\bigwedge^qM_1^t$. Proper base change applies to the compact
Hausdorff fibres
\cite[Tag 09V5, Lemma 20.18.1; chapter PDF p.~30]{StacksProject}.
Writing $R^q=R^q\pi_*\mathbb Z$ for the constructed map, the exact
integral image and primitive-minor certificates in
Theorem~\ref{maps:src:F.2.6} therefore give
\[
\begin{array}{c|ccc}
q&H^0(\mathbb P^1,R^q)&H^1(\mathbb P^1,R^q)&H^2(\mathbb P^1,R^q)\\\hline
1&\mathbb Z&0&\mathbb Z/2\\
2&\mathbb Z/2&0&\mathbb Z/2\\
3&\mathbb Z/2&\mathbb Z^3&\mathbb Z/2\\
4&\mathbb Z\oplus\mathbb Z/2&0&\mathbb Z/2.
\end{array}
\]
All fibres are connected, so $R^0=\mathbb Z$. In the integral Leray
sequence
\cite[Tag 01F2, Lemma 20.13.4; chapter PDF p.~22]{StacksProject},
the three maps
\[
 d_2^{0,2}:\mathbb Z/2\longrightarrow\mathbb Z/2,
 \qquad
 d_2^{0,3}:\mathbb Z/2\longrightarrow\mathbb Z/2,
 \qquad
 d_2^{0,4}|_{\rm tors}:\mathbb Z/2\longrightarrow\mathbb Z/2
\]
are the isomorphisms identified by the actual finite boundary maps
in Step~2. The base has cohomological dimension two, so there are no
later differentials.

\item Theorem~\ref{night:P1.10} gives $\pi_1(M_c)=0$, hence
$H_1(M_c;\mathbb Z)=H^1(M_c;\mathbb Z)=0$. The total-degree-one term
in the Leray sequence is precisely the kernel of
\[
 d_2^{0,1}:H^0(R^1)=\mathbb Z
                 \longrightarrow H^2(R^0)=\mathbb Z.
\]
Thus this map is multiplication by a nonzero integer. In total
degree two, Step~3 kills $E_2^{0,2}$ and $E_2^{1,1}=0$; consequently
$H^2(M_c;\mathbb Z)$ is the finite cokernel of $d_2^{0,1}$.
The universal coefficient sequence, with $H_1=0$, identifies it with
$\operatorname{Hom}(H_2(M_c;\mathbb Z),\mathbb Z)$
\cite[Theorem 3.2, p.~195, and Corollary 3.3, p.~196]{Hatcher2002}.
This is a free group, since the homology of the compact smooth
manifold is finitely generated. Its finiteness forces $H^2=0$ and
forces $d_2^{0,1}$ to be a unit.

\item Step~3 kills every total-degree-three term, so $H^3=0$.
In total degree four the surviving terms are $\mathbb Z^3$ at
$(1,3)$ and the kernel of $d_2^{0,4}$ at $(0,4)$. Choose a free
generator of $\mathbb Z\oplus\mathbb Z/2$ and let
$\epsilon\in\mathbb Z/2$ be its image. Then
\[
 d_2^{0,4}(n,t)=\epsilon n+t,
 \qquad
 \ker d_2^{0,4}=\mathbb Z(1,\epsilon).
\]
The Leray filtration therefore gives an extension of $\mathbb Z$
by $\mathbb Z^3$. It splits because the quotient is free, and
$H^4(M_c;\mathbb Z)=\mathbb Z^4$.

The universal coefficient sequence with $H^3=0$ gives
$\operatorname{Ext}(H_2,\mathbb Z)=0$ and
$\operatorname{Hom}(H_3,\mathbb Z)=0$. Thus $H_2$ is free and $H_3$
is finite. Since $H^2=\operatorname{Hom}(H_2,\mathbb Z)=0$, we have
$H_2=0$. The next universal coefficient sequence injects the finite
group $\operatorname{Ext}(H_3,\mathbb Z)$ into the free group $H^4$;
it is zero, and hence $H_3=0$.

The complex orientation and integral Poincar\'e duality
\cite[Theorem 3.30, p.~241]{Hatcher2002} give
\[
 H_4=H^4=\mathbb Z^4,\qquad
 H_5=H^3=0,\quad H_6=H^2=0,\quad H_7=H^1=0,\quad H_8=H^0=\mathbb Z.
\]
Connectedness gives $H_0=\mathbb Z$. Finally the Hurewicz maps
\cite[Theorem 4.32, p.~366]{Hatcher2002}, applied successively to the
simply connected manifold, give $\pi_2=H_2=0$ and then
$\pi_3=H_3=0$.
\end{enumerate}
\end{proof}
